% Part: incompleteness
% Chapter: incompleteness-provability
% Section: provability-conditions

\documentclass[../../../include/open-logic-section]{subfiles}

\begin{document}

\olfileid{inc}{inp}{prc}

\olsection{$\Th{PA}$ साठीच्या \usetoken{S}{derivability} अटी}

पेआनो अंकगणित, म्हणजे $\Th{PA}$, ही
सर्व !!{formula}s साठी विगमनाची स्वयंसिद्धके
जोडून मिळणारी $\Th{Q}$ ची विस्तार-उपपत्ती
आहे. म्हणजे प्रत्येक !!{formula}~$!A$
साठी $\Th{Q}$ मध्ये पुढील रूपाचे
स्वयंसिद्धक जोडतात:
\[
(!A(0) \land \lforall[x][(!A(x) \lif !A(x'))]) \lif \lforall[x][!A(x)]
\]
या रूपबंधात इतर मुक्त चल असतील, तर आरंभी त्यांच्यावर
सार्विक संख्यापक जोडून वाक्यरूपाचे स्वयंसिद्धक घेतो.
हे खरे तर एक \emph{रूपबंध} आहे, म्हणजे
अनंत स्वयंसिद्धके आहेत, हे लक्षात घ्या
(आणि $\Th{PA}$ चे सांत
स्वयंसिद्धकसंचाने स्वयंसिद्धकीकरण करता येत
{\em नाही} असेही सिद्ध होते).
पण चिन्हमाला विगमनाच्या
स्वयंसिद्धकाचे उदाहरण आहे
की नाही हे प्रभावीपणे ठरवता
येत असल्यामुळे $\Th{PA}$
चा स्वयंसिद्धकसंच संगणनक्षम
आहे. $\Th{PA}$ ही~$\Th{Q}$
पेक्षा बरीच अधिक सबल
उपपत्ती आहे. उदाहरणार्थ,
नेहमीच्या पद्धतीने विगमन
वापरून बेरीज आणि गुणाकार
क्रमनिरपेक्ष आहेत हे
सहज सिद्ध करता येते.
खरे तर बहुतेक सांत
संख्यासिद्धांतिक आणि
संयोजनात्मक युक्तिवाद
$\Th{PA}$ मध्ये मांडता येतात.

$\Th{PA}$ ही संगणनक्षमपणे
स्वयंसिद्धकीकृत असल्यामुळे
$\Prf[\Th{PA}](x,y)$ हे
!!{derivability} विधेय
संगणनक्षम आहे आणि म्हणून
$\Th{Q}$ मध्ये निरूपित
होते (म्हणूनच
$\Th{PA}$ मध्येही).
पूर्वीप्रमाणे, या संबंधाचे
निरूपण करणाऱ्या सूत्रासाठी
$\OPrf[\Th{PA}](x,y)$
वापरू. $\OProv[\Th{PA}](y)$
हे सूत्र
$\lexists[x][\OPrf[\Th{PA}](x,y)]$
असे मानू. त्याचा सहज
अर्थ असा: “$y$ हे
$\Th{PA}$ च्या स्वयंसिद्धकांपासून
!!{derivable} आहे.”
$\Th{Q}$ च्या स्वयंसिद्धकांपेक्षा
थोडे अधिक सामर्थ्य लागते,
कारण आपण वापरत असलेली
उपपत्ती या !!{derivability}
विधेयाविषयीची काही मूलभूत
तथ्ये !!{derive} करू शकली
पाहिजे. पुढील अटींसाठी प्रत्यक्ष आकारिक सिद्धतांच्या प्रमाण
संकेतांकीकरणावर आधारित विधेय निवडले आहे. केवळ निर्णेय
संबंधाचे निरूपण करणारे मनमानी सूत्र घेतल्याने या अटी
आपोआप मिळत नाहीत. नेमकी आवश्यक तथ्ये
पुढीलप्रमाणे:
\begin{enumerate}
\item[P1.] जर $\Th{PA} \Proves !A$, तर $\Th{PA} \Proves
  \OProv[\Th{PA}](\gn{!A})$.
\item[P2.] सर्व !!{formula}s $!A$ आणि $!B$ साठी,
  \[
  \Th{PA} \Proves \OProv[\Th{PA}](\gn{!A \lif !B}) \lif
  (\OProv[\Th{PA}](\gn{!A}) \lif \OProv[\Th{PA}](\gn{!B})).
  \]
\item[P3.] प्रत्येक !!{formula}~$!A$ साठी,
  \[
  \Th{PA} \Proves \OProv[\Th{PA}](\gn{!A})
  \lif \OProv[\Th{PA}](\gn{\OProv[\Th{PA}](\gn{!A})}).
  \]
\end{enumerate}
हे तीन गुणधर्म खरे आहेत
याची खात्री करण्यासाठी
$\OProv[\Th{PA}](y)$
या !!{formula} चे बारकाईने
वर्णन करून संबंधित आकारिक
!!{derivation}s चे वर्णन
$\Th{PA}$ च्या स्वयंसिद्धकांनी
करावे लागते. अटी (1) आणि~(2)
सोप्या आहेत; खरे काम
अट~(3) साठी लागते
(त्यासाठी नेमके काय
करावे लागेल याचा विचार करा
\dots). तपशील मांडणे
कंटाळवाणे आणि फारसे
रंजक नसल्यामुळे
$\Th{PA}$ मध्ये वरचे
तीन गुणधर्म आहेत हे
येथे आपण गृहीत धरू.
$\OProv[\Th{PA}](y)$
ची योग्य निवड पुढील
अटही पूर्ण करेल:
\begin{enumerate}
\item[P4.] जर $\Th{PA} \Proves \OProv[\Th{PA}](\gn{!A})$, तर
  $\Th{PA} \Proves !A$.
\end{enumerate}
पण हे तथ्य आपल्याला
लागणार नाही. ही शेवटची अट आधीच्या तीन आकारिक अटींचा
निष्कर्ष नाही; तिच्यासाठी पेआनो अंकगणिताचे मानक प्रतिमानातील
सत्य आणि प्रत्यक्ष सिद्धतासंकेतांकांचे प्रमाण विधेय वापरले आहे.

\begin{digress}
तपशील वगळण्याबाबत
गोडेलनेही आपल्यासारखाच
आळस केला. 1931 च्या
शोधपत्राच्या अखेरीस
त्याने दुसऱ्या अपूर्णता
प्रमेयाच्या सिद्धतेची
रूपरेषा दिली आणि
तपशील पुढील शोधपत्रात
देण्याचे आश्वासन दिले.
ते तपशील त्याने
नंतर प्रकाशित केले
नाहीत; युक्तिवाद समजणाऱ्या
सर्वांनाच तो पूर्ण करता
येईल असे वाटत असल्यामुळे
तपशील भरण्याची गरजही
त्याला पडली नाही.
\end{digress}

\end{document}
